\documentclass{article}
\usepackage{lmodern}
\usepackage{amsmath}
\usepackage{listings}
\usepackage{fullpage}
\usepackage{parskip}
\usepackage{xcolor}
\lstset{
	basicstyle=\ttfamily,
	columns=fullflexible,
	frame=single,
	breaklines=true,
	postbreak=\mbox{\textcolor{red}{$\hookrightarrow$}\space},
}
\begin{document}
\title{Experiment `p\_3\_4\_classification' Results}
\author{\today}
\date{}
\maketitle
\textbf{Experiment outcome: }SUCCESS
\\\textbf{Bad responses: }0
\\\textbf{Responses containing}~\texttt{assume}~\textbf{: }0
\\\textbf{Responses containing}~\texttt{expect}~\textbf{: }0
\\\textbf{Resolution attempts: }1
\\\textbf{Hard fails (resolution): }0
\\\textbf{Soft fails (resolution): }0
\\\textbf{Verification attempts: }1
\section*{Problem Specification}
\textbf{Problem name: }p\_3\_4\_classification
\\\textbf{Natural language statement: }null
\\\textbf{Method signature: }p\_3\_4\_classification (i: int, j: int, k: int) returns (s: string)
\subsection*{Ensures}
\begin{itemize}
\item \texttt{if i == j == k then s == "all the same" else if i != j \&\& j != k \&\& i != k then s == "all different" else s == "neither"}
\end{itemize}
\clearpage
\section*{GenAI interactions}
Below you will find all interactions between the `user' (program) and the `assistant' (OpenAI).
\subsection*{Program $\rightarrow$ GenAI}
\begin{lstlisting}
You are given the following task to perform in Dafny:

Write a method to satisfy the below requirements.

The signature should be:

p_3_4_classification (i: int, j: int, k: int) returns (s: string)

The method should respect the following contract:

ensures if i == j == k then s == "all the same" else if i != j && j != k && i != k then s == "all different" else s == "neither"

Produce and show only the Dafny body of this method, including the curly braces that surround it. Do not show the signature nor contract. You must not use 'assume' nor 'expect' anywhere in your code.
\end{lstlisting}
\subsection*{GenAI $\rightarrow$ Program}
\textbf{System fingerprint: }null
\\\textbf{ID: }chatcmpl-Dn39UqXhhBzHnbmHdNUf8iXLv8dxC
\\\textbf{Model: }gpt-5.5-2026-04-23
\\\textbf{Created at: }1780582560
\\\textbf{Finish reason: }stop
\\\textbf{Usage: }332 tokens in, and 200 tokens out
\begin{lstlisting}
{
  if i == j && j == k {
    s := "all the same";
  } else if i != j && j != k && i != k {
    s := "all different";
  } else {
    s := "neither";
  }
}
\end{lstlisting}
\section*{Final Program}
\begin{lstlisting}
method p_3_4_classification (i: int, j: int, k: int) returns (s: string)
	ensures if i == j == k then s == "all the same" else if i != j && j != k && i != k then s == "all different" else s == "neither"
{
  if i == j && j == k {
    s := "all the same";
  } else if i != j && j != k && i != k {
    s := "all different";
  } else {
    s := "neither";
  }
}
\end{lstlisting}
\section*{Total Token Usage}
\textbf{Input tokens: }332
\\\textbf{Output tokens: }200
\\\textbf{Reasoning tokens: }116
\\\textbf{Sum of `total tokens': }532
\section*{Experiment Timings}
\textbf{Overall Experiment} started at 1780582559727, ended at 1780582595512, lasting 35785ms (35.79 seconds)
\\\textbf{Iteration \#1} started at 1780582559729, ended at 1780582595512, lasting 35783ms (35.78 seconds)
\end{document}
